% ============================================================================
% Two electrons on helium as a differential sensor for Majorana zero modes
% Target: Phys. Rev. A / PRX Quantum (adjust class options below once decided)
% ============================================================================
\documentclass[aps,pra,reprint,superscriptaddress,longbibliography]{revtex4-2}
% \documentclass[aps,prxquantum,reprint,superscriptaddress,longbibliography]{revtex4-2}

\usepackage{graphicx}
\usepackage{amsmath,amssymb}
\usepackage{bm}
\usepackage{hyperref}

\begin{document}

\title{Two Electrons on Helium as a Differential Sensor for Majorana Zero Modes}
% Working title — revisit once the specific signal channel (Sec.~III) is chosen.

\author{Author list TBD}
\affiliation{Affiliation TBD}

\date{\today}

\begin{abstract}
We propose using a pair of Coulomb-entangled electrons on superfluid helium, prepared in the
motional/spin singlet configuration used as a differential field sensor in
Ref.~\cite{perruzza2026}, to probe local signatures associated with Majorana zero modes (MZMs) in
a proximal topological superconductor. Because the singlet is insensitive to common-mode
perturbations and maximally sensitive to a \emph{difference} between the two trap sites, positioning
the two electrons asymmetrically with respect to a topological wire converts a local,
sub-shot-noise MZM signature into a differential signal read out through the same
quantum-Fisher-information framework already developed for particle-physics sensing on this
platform. This is a proposal: we outline the signal channels under consideration, the required
theoretical extension of the existing singlet--triplet sensor model, and the open questions that
must be resolved before a concrete protocol can be simulated.
\end{abstract}

\maketitle

\section{Introduction}

Two electrons on superfluid helium, entangled through their unscreened Coulomb interaction, have
been shown theoretically to support a differential-sensing scheme: prepared in the motional/spin
singlet, the pair is blind to common-mode fields but maximally responsive to a \emph{gradient}
between the two trap sites, quantified through the quantum and classical Fisher information and
benchmarked against an independent-spin Ramsey scheme~\cite{perruzza2026}. That proposal targeted
generic magnetic-field gradients (e.g.\ from light dark matter or axion-like fields). The
underlying entangled states, and the machinery for computing them (Hartree/CI on a sinc-DVR basis,
configurations I/II/III), were established for the two-electron helium platform in
Ref.~\cite{beysengulov2024}, with time-dependent control of the same states realized as high-fidelity
two-qubit gates in Ref.~\cite{leinonen2025}. Reference~\cite{perruzza2026} also grounds the sensing
readout of a real device in Refs.~\cite{koolstra2025impedance,castoria2025}, which demonstrate the
resonator-based charge/orbital readout hardware this platform actually has today.

Majorana zero modes (MZMs), predicted to occur at the ends of one-dimensional topological
superconductors~\cite{kitaev2001unpaired} and in proximitized semiconductor
nanowires~\cite{lutchyn2010majorana,oreg2010helical}, have been pursued experimentally for over a
decade, with zero-bias conductance signatures reported in hybrid
superconductor--semiconductor nanowires~\cite{mourik2012signatures} and reviewed extensively
since~\cite{beenakker2013search,lutchyn2018majorana}. A central experimental difficulty is that MZM
signatures are subtle, easily mimicked by trivial (non-topological) Andreev bound states, and
sensitive to quasiparticle poisoning — a stray quasiparticle tunneling onto the topological segment
flips its fermion parity and can toggle any locally coupled signal. This is precisely the kind of
weak, potentially transient, spatially localized perturbation that a differential entangled-pair
sensor is suited to pick up while rejecting the (much larger) common-mode electromagnetic
environment of a cryostat.

This project asks whether the existing eHe differential-sensor machinery can be repurposed, with
the ``target signal'' reinterpreted from a magnetic-field gradient to whatever local field or charge
signature a nearby MZM produces. We do not yet have a settled answer to \emph{which} physical
channel (Sec.~\ref{sec:model}) is both (i) strong enough to be detectable at achievable electron--source
distances and (ii) compatible with the existing singlet--triplet coupling operator. The goal of this
manuscript, at this stage, is to lay out the extension of the model and the open questions rather
than to report a result.

\section{Model}
\label{sec:model}

\subsection{Reusing the existing sensor Hamiltonian}

The starting point is the effective two-level (ground/excited, singlet-like/triplet-like) model of
Ref.~\cite{perruzza2026}:
\begin{equation}
H_{\rm eff} = E_g|g\rangle\langle g| + E_e|e\rangle\langle e| +
\lambda(t)\left(|g\rangle\langle e| + |e\rangle\langle g|\right),
\end{equation}
where in the original proposal $\lambda = (g\mu_B \delta B/2) M_{ge}$ couples to a magnetic-field
gradient $\delta B = B_L - B_R$ via the matrix element $M_{ge} = \langle e|S^z_L - S^z_R|g\rangle$.
For a Majorana-mode signal, $\lambda(t)$ should instead encode whatever locally asymmetric quantity
the MZM produces at the two trap sites — to be determined (see below). The quantum/classical Fisher
information and Cram\'er--Rao bound machinery of Ref.~\cite{perruzza2026} carries over unchanged
once $\lambda(t)$ is specified.

\subsection{Candidate signal channels (to be evaluated)}

\begin{itemize}
  \item \textbf{Parity-switching noise.} A quasiparticle-poisoning event on the topological segment
  flips the ground-state fermion parity, which (depending on the specific hybrid architecture) can
  modulate a local charge or field asymmetric between the two trap sites. This would appear as a
  telegraph-noise-like signal rather than a coherent Rabi oscillation, requiring a different Fisher
  information treatment (parameter = switching rate, not a static/slowly-varying field).
  \item \textbf{Stray field from Majorana-induced edge/Andreev currents.} If the topological wire
  supports a weak stray magnetic field at its end (from the associated Andreev bound state charge
  distribution or edge current), this maps most directly onto the existing $\delta B$ channel and
  requires the smallest change to the Ref.~\cite{perruzza2026} formalism.
  \item \textbf{Direct charge-coupling channel.} Reusing the electric susceptibility / dispersive
  framework of Ref.~\cite{castoria2025} (an electron's motional state already couples to local
  charge/potential fluctuations through the resonator), a nearby MZM's contribution to the local
  electrostatic potential could couple directly to $(n_L - n_R)$ rather than to spin — a fundamentally
  different (charge, not spin-gradient) sensing channel that would need its own generator and its own
  QFI expression.
\end{itemize}

Each channel implies a different coupling operator, a different noise model, and a different
distance/coupling-strength budget; identifying the most promising one (or ruling all three out at
currently achievable electron--nanowire separations) is the first concrete task for this project.

\subsection{Open theoretical questions}

\begin{itemize}
  \item What is the realistic magnitude of the candidate signal at achievable electron--MZM
  separations, and how does it compare to the $\delta B$ magnitudes assumed in
  Ref.~\cite{perruzza2026}?
  \item Is the relevant noise process (e.g.\ quasiparticle poisoning) better modeled as a static
  parameter to estimate (Fisher information / Cram\'er--Rao, as in the existing framework) or as a
  stochastic process requiring a different estimation-theoretic treatment (e.g.\ change-point
  detection, quantum non-demolition parity readout)?
  \item Does the factor-of-2 QFI advantage carry over quantitatively, or does the new coupling
  operator's matrix element $M_{ge}$ (or its charge-channel analogue) reduce the achievable gain?
  \item What does the required electron--nanowire integration look like physically (shared substrate?
  proximity-coupled chip?), and is it compatible with the trap/reservoir geometry of
  Ref.~\cite{castoria2025}?
\end{itemize}

\section{Results}

Not yet available — this section will report the candidate-channel comparison and, for the
strongest candidate, a QFI/CRB estimate analogous to Sec.~IV of Ref.~\cite{perruzza2026}.

\section{Discussion and outlook}

If a viable signal channel is identified, next steps are: (i) extend the CI/Hartree machinery of
Ref.~\cite{beysengulov2024} to include the new coupling operator, reusing the numerical toolkit
described in the project's \texttt{ci-dvr-hartree-numerics} skill; (ii) simulate the achievable QFI
as a function of electron--MZM distance and compare against the independent-spin (Ramsey) benchmark;
(iii) assess whether RL-based trap optimization (see the companion \texttt{rl-agentic-quantum-control}
project) can improve the achievable coupling by reshaping the trap potential, as it already does for
the field-gradient sensor.

\begin{acknowledgments}
Placeholder.
\end{acknowledgments}

\bibliography{references}

\end{document}
